\section{The Scent of a Woman}

This is why it is important to read older books, since it is easy to forget how quickly customs change, at least for those inhabiting the far West. In Gustav Meyer's The Green Face, we find this description of the seductive salesgirl in the magic shop:

\begin{quotex}
As she did so demonstrate a magic trick for him she brought the whole range of a practised feminine charm into play, from the breasts deliberately displayed to the male, to the discreet, almost telepathic scent given off by her skin \emph{which she intensified by occasionally lifting her arm to send a supplementary blast from the armpit}.
\end{quotex}


It certainly worked on the patron, since he bought the trick even though he neither wanted it nor needed it. Men today, who expect women to smell like a Dove beauty bar\footnote{\url{https://www.dove.com/us/en/washing-and-bathing/beauty-bar/white-beauty-bar.html}}, probably won't understand how her armpit scent is an attraction. Or perhaps they forget the natural connection between body odor and subconscious human sexual attraction\footnote{\url{https://en.wikipedia.org/wiki/Body\_odor\_and\_subconscious\_human\_sexual\_attraction}}. Or perhaps she was a sorceress practicing her charm ... 

\flrightit{Posted on 2012-08-23 by Cologero}

\begin{center}* * *\end{center}

\begin{footnotesize}
\begin{sffamily}
\texttt{Aghorable on 2012-08-24 at 01:47 said:}

Of course. This is why some women who might not be considered sexy if judged by popular standards still get a man's juices flowing when he is in their presence.

Oh, and Charlotte obviously could not resist returning after you flashed that deadly handlebar `stache and white suit.

\hfill

\texttt{Cologero on 2012-08-24 at 16:55 said:}

I am only into chapter 3, but I see Meyrink as a colorful writer, with the ability to describe various human types in great detail. He punctuates the story with ruminations on destiny and the intricacies of the thinking process. I can understand why Evola would have been attracted to this work. Meyrink describes the hypocrisy of Christian bourgeois life in Amsterdam in the ruins of post-WWI. For example, he describes one woman, a ``humanitarian'', who runs charitable events for people in the abstract, but does nothing for real individual persons. She has all the politically correct thoughts, although they do not go very deep. I was surprised, too, to learn that the Nigerian email scam had its own version at that time (without the Nigerian). I don't know yet why Guenon thought it was such a evil book, perhaps that will come out later on.

\hfill

\texttt{Logres on 2012-08-24 at 20:59 said:}

I think it's good you are ``testing the spirits'' by reading the book yourself -- I have one book by M that I could read, and am interested in Guenon's perspective, given he was so critical of so many of the figures Evola was interested in.

\hfill

\texttt{Logres on 2012-08-24 at 20:59 said:}

Not that you have to ``test everything'', of course.

\hfill

\texttt{Cologero on 2012-08-24 at 21:13 said:}

Actually, Logres, I'm reading it for pleasure. Even Eliade enjoyed Meyrink. Guenon, in his letters, claims that presenting Traditional information as a ``parody'' is an indication of a counter-initiation. We'll see ... or not.

BTW, oddly enough Guenon does recommend reading that type of material or, as you put it, ``test everything''.

\hfill

\texttt{Cologero on 2012-08-26 at 18:53 said:}

I am still far from the conclusion that Guenon objects to. There are a series of macabre events and eccentric characters that form the story. However, there are some serious monologues about the nature of thought and destiny that, I believe, few men could write, and probably not too many more could even understand.

\hfill

\texttt{Cologero on 2012-08-28 at 11:41 said:}

From The Green Face

\begin{quote}\footnotesize

Until now men have torn each other apart for the sake of certain dubious invisible beings that are careful not to call themselves spirits, but `ideals'. I think the hour has finally come for the war against these invisible enemies, and I would like to play a part in it. For years I have been aware that I was being trained in spiritual warfare, but until now I have never had this clear feeling that a great battle against these damned ghosts is at hand. I tell you, once you start clearing out all those false ideals, there's no end to it. \emph{You've no idea what piles of humbug brazenly posing as truth you have inherited.}
\end{quote}


\hfill

\texttt{Numen on 2017-08-24 at 20:51 said:}

I've recently come a across a books titled ``Scenting Salvation: Ancient Christianity and the Olfactory Imagination,'' which is about the aspects of scent and smell in early Christianity and Western mysticism. I have yet to fully read it in it's entirety, but some interesting selections just from skimming through:

``exuding sweet fragrance to disguise the bad odor of jealousy'' -- Palladius

``Dearer to God are the trials for righteousness' sake than all vows and sacrifices. And dearer too, is the odor of the sweat of fatigue they cause, than all the drugs of sweet scent and exquisite perfumes'' -- Isaac of Nineveh

``virgins ought to refrain from using too much perfume''


\hfill

\end{sffamily}
\end{footnotesize}

